\documentclass[UTF8,a4paper,12pt]{ctexart}

\usepackage{geometry}
\geometry{left=2.2cm,right=2.2cm,top=2.2cm,bottom=2.2cm}

\usepackage{array}
\usepackage{longtable}
\usepackage{booktabs}
\usepackage{hyperref}
\usepackage{xcolor}
\usepackage{ragged2e}
\usepackage{enumitem}

\hypersetup{
    colorlinks=true,
    linkcolor=black,
    urlcolor=black
}

\setlength{\tabcolsep}{3pt}
\renewcommand{\arraystretch}{1.35}
\sloppy

\newcolumntype{P}[1]{>{\RaggedRight\arraybackslash}p{#1}}
\newcommand{\pass}{通过}

\title{\textbf{AI学术鉴伪系统测试用例报告}}
\author{刘宇童 \space 聂薪羽 \space 苏天 \space 翟润 \space 张与耕}
\date{2026年5月23日}

\begin{document}

\maketitle

\section*{测试说明}

本报告依据学术内容诚信检测平台最终测试结果整理，采用“任务一（手动测试）”与“任务二（脚本测试）”两部分组织测试用例。任务一面向用户端与管理端核心业务流程，共43项功能测试用例；任务二面向后端接口与边界条件，共30项接口测试用例。两类测试均已执行完成，测试结果均为“通过”。

测试环境为本地联调环境：AI服务端口为8010，Django业务API端口为8000，用户端端口为3000，管理端端口为3001。测试覆盖编辑、专家、组织管理员、软件管理员等角色，重点验证学术检测、检测历史、综合鉴伪报告、人工审核闭环、管理端运维以及后端接口边界校验等核心功能。

\section{任务一（手动测试）}

针对学术鉴伪系统用户端与管理端的核心业务流程进行功能性需求测试。任务一为人工手动测试，测试人员按步骤在前端页面执行操作并观察页面反馈。根据最终测试报告，任务一共43项功能测试用例，实际执行结果均为“通过”。

\begin{longtable}{|P{0.9cm}|P{3.5cm}|P{4.7cm}|P{5.0cm}|P{1.2cm}|}
\hline
\textbf{测试序号} & \textbf{前置条件} & \textbf{步骤} & \textbf{期望结果} & \textbf{测试结果} \\
\hline
\endfirsthead

\hline
\textbf{测试序号} & \textbf{前置条件} & \textbf{步骤} & \textbf{期望结果} & \textbf{测试结果} \\
\hline
\endhead

1 & 系统已启动，用户端 \url{http://127.0.0.1:3000} 可访问，存在编辑用户账号。 & 进入用户端登录页，输入编辑用户邮箱、密码和验证码，选择编辑身份登录。 & 登录成功，进入主页；侧栏包含学术检测、检测历史、人工审核申请等编辑侧入口。 & \pass \\
\hline

2 & 系统已启动，存在专家用户账号。 & 进入用户端登录页，输入专家用户邮箱、密码和验证码，选择专家身份登录，并尝试访问编辑侧上传路径。 & 登录成功后仅显示人工审核、个人主页等专家侧入口；访问 \url{/upload} 时重定向至 \url{/review}。 & \pass \\
\hline

3 & 系统已启动，管理端 \url{http://127.0.0.1:3001} 可访问，存在组织管理员账号。 & 进入管理端登录页，输入组织管理员账号、密码和验证码登录。 & 登录成功后进入管理端工作台；侧栏包含审批、系统日志等管理端入口。 & \pass \\
\hline

4 & 用户端登录页可访问。 & 在登录页输入错误密码或错误验证码后点击登录。 & 登录失败，页面给出错误提示；不进入系统，不保存有效登录状态。 & \pass \\
\hline

5 & 编辑用户已登录。 & 进入个人主页查看账号信息。 & 页面展示用户信息；发布者可见等级与默认检测模式等信息。 & \pass \\
\hline

6 & 用户已登录。 & 点击退出登录，并再次访问需要认证的页面。 & 系统退出登录并跳转至登录页；再次访问受保护页面时要求重新登录。 & \pass \\
\hline

7 & 编辑用户已登录，准备合法图像文件。 & 进入上传/检测页面，选择图像检测并选择图片文件。 & 页面显示文件名与文件类型，文件进入待检测队列，可开始检测。 & \pass \\
\hline

8 & 编辑用户已登录，准备不支持的文件类型，例如 \texttt{.exe}。 & 在图像检测上传区域选择不支持的文件。 & 系统拒绝上传或给出格式不合法提示，不形成有效检测任务。 & \pass \\
\hline

9 & 编辑用户已登录，已成功上传合法图片，AI服务在线。 & 选择标准模式并提交图像AI检测任务。 & 检测历史新增任务，任务状态最终为已完成。 & \pass \\
\hline

10 & 编辑用户已登录，已成功上传合法图片，AI服务在线。 & 选择精准模式并提交图像AI检测任务。 & 任务提交成功并完成检测，检测记录能够体现或区分精准模式。 & \pass \\
\hline

11 & 编辑用户已登录，存在已提交的图像检测任务。 & 进入检测历史页面，查看任务列表并打开任务详情。 & 详情页展示任务信息、检测状态、报告入口；URL中包含 \texttt{detail\_id}。 & \pass \\
\hline

12 & 编辑用户已登录，检测历史中存在多种类型或状态的任务。 & 在检测历史页面按状态、类型或时间进行筛选。 & 列表仅展示符合筛选条件的任务；清空筛选后恢复默认列表。 & \pass \\
\hline

13 & 编辑用户已登录，准备正常PDF论文文件。 & 进入上传/检测页面，选择论文检测并上传PDF文件。 & PDF文件上传成功，可进入后续论文检测流程。 & \pass \\
\hline

14 & 编辑用户已登录，准备非PDF文件。 & 在论文检测入口上传非PDF文件。 & 系统拒绝上传，提示仅支持PDF文件，不创建论文检测任务。 & \pass \\
\hline

15 & 编辑用户已登录，论文PDF已上传成功。 & 提交全篇论文AIGC检测。 & 系统创建论文AIGC检测任务，检测历史中出现对应记录并最终完成。 & \pass \\
\hline

16 & 编辑用户已登录，论文PDF已上传成功。 & 提交论文资源检查任务。 & 系统创建资源检查任务，检测历史中可查看资源检查记录。 & \pass \\
\hline

17 & 编辑用户已登录，准备正常Review文本。 & 进入Review检测入口，粘贴Review文本并提交。 & 系统完成文本清洗并创建Review检测任务，页面显示提交成功和任务状态。 & \pass \\
\hline

18 & 编辑用户已登录，Review检测入口可访问。 & 不输入文本且不上传 \texttt{.txt} 文件，直接点击提交。 & 系统拒绝提交，提示Review内容不能为空或需上传有效文件。 & \pass \\
\hline

19 & 编辑用户已登录，准备 \texttt{.txt} 格式Review文件。 & 选择 \texttt{.txt} 文件并提交Review检测。 & 文件上传成功，系统创建Review检测任务，检测历史中显示Review任务。 & \pass \\
\hline

20 & 编辑用户已登录，准备包含图片的ZIP压缩包。 & 在批量提交入口上传ZIP压缩包并开始检测。 & 系统从ZIP中抽取图片并进入检测流程，检测历史中出现批次任务。 & \pass \\
\hline

21 & 编辑用户已登录，准备DOCX文件。 & 在批量提交入口尝试提交DOCX文件。 & 系统提示需转换为PDF，不产生论文检测任务。 & \pass \\
\hline

22 & 编辑用户已登录，存在已完成的论文AIGC检测任务。 & 打开论文AIGC检测结果页。 & 页面展示段落风险、事实性鉴伪、摘录高亮等结果信息。 & \pass \\
\hline

23 & 编辑用户已登录，系统存在可汇总的多模态检测结果。 & 进入综合鉴伪报告页面查看报告。 & 报告包含各模态小节与综合结论，内容完整可读。 & \pass \\
\hline

24 & 编辑用户已登录，综合鉴伪报告已生成。 & 点击下载综合鉴伪PDF报告。 & 浏览器开始下载PDF文件，下载后可正常打开。 & \pass \\
\hline

25 & 编辑用户已登录，综合鉴伪报告已生成。 & 点击导出综合鉴伪HTML报告。 & 浏览器下载包含结构化内容的HTML文件。 & \pass \\
\hline

26 & 编辑用户已登录，存在已完成或可查看的检测任务。 & 在检测详情或人工审核页面填写审核理由并发起人工审核申请。 & 申请提交成功，状态显示为待管理员审批；申请列表中出现该记录。 & \pass \\
\hline

27 & 编辑用户已登录，存在已提交的人工审核申请。 & 进入人工审核申请列表，按状态或时间筛选申请记录。 & 列表筛选结果正确，可查看申请状态、提交时间、任务类型和申请理由。 & \pass \\
\hline

28 & 组织管理员已登录管理端，存在待审批人工审核申请。 & 进入管理端人工审核页面，打开待审批申请详情。 & 详情页展示申请人、检测任务、申请理由、任务类型、提交时间等信息。 & \pass \\
\hline

29 & 组织管理员已登录，存在待审批人工审核申请。 & 在人工审核详情中点击通过申请。 & 申请状态更新为已通过；专家侧任务池出现对应人工审核任务。 & \pass \\
\hline

30 & 组织管理员已登录，存在另一条待审批人工审核申请。 & 在人工审核详情中选择拒绝并填写拒绝原因。 & 申请状态更新为已拒绝；编辑侧可查看拒绝原因；专家侧不会出现该任务。 & \pass \\
\hline

31 & 专家用户已登录，管理员已通过至少一条人工审核申请。 & 进入专家审阅页面查看任务池。 & 任务池展示分配给当前专家的待处理任务，可看到任务类型、发布者、状态和提交时间。 & \pass \\
\hline

32 & 专家用户已登录，任务池存在图像类人工审核任务。 & 点击图像类人工审核任务进入详情页。 & 详情页展示待审核图片、AI检测参考结果、评分或结论填写区域和提交按钮。 & \pass \\
\hline

33 & 专家用户已登录，任务池存在论文或Review类人工审核任务。 & 点击论文或Review类人工审核任务进入详情页。 & 详情页展示文本片段、AI摘要或风险点、人工结论填写区域和提交按钮。 & \pass \\
\hline

34 & 专家用户已进入人工审核详情页。 & 填写审核结论和理由后提交人工审核结果。 & 提交成功，任务状态更新为已完成；再次进入列表时显示完成状态。 & \pass \\
\hline

35 & 编辑用户已登录，专家已完成某条人工审核任务。 & 进入人工审核结果汇总页面查看该任务。 & 页面展示管理员审批状态、专家审核状态、专家结论、说明和综合汇总信息。 & \pass \\
\hline

36 & 编辑用户已登录，存在已完成检测任务。 & 在检测结果详情页点击下载单任务检测报告。 & 已完成任务可下载PDF报告；未完成任务给出明确提示。 & \pass \\
\hline

37 & 管理员已登录管理端，系统存在多条检测任务。 & 进入检测任务页面，按任务类型、状态或用户筛选，并打开任务详情。 & 任务列表按筛选条件正确刷新，任务详情与检测状态显示正常。 & \pass \\
\hline

38 & 管理员已登录管理端，系统存在文件资源。 & 进入资源与文件页面，按文件类型、上传用户或时间筛选。 & 文件列表与筛选功能正常，可查看文件基本信息与关联任务。 & \pass \\
\hline

39 & 管理员已登录管理端，系统存在用户和组织数据。 & 进入用户与组织页面，搜索某一用户名或组织名。 & 搜索结果正确过滤，用户角色、所属组织、权限等信息显示正确。 & \pass \\
\hline

40 & 管理员已登录管理端，系统存在操作日志和检测任务记录。 & 进入日志页面，按操作类型或时间范围筛选日志。 & 操作留痕与检测任务记录可查，日志筛选结果正确。 & \pass \\
\hline

41 & 发布者用户已登录。 & 查看用户端侧栏入口。 & 侧栏仅保留一个“学术检测”入口，无重复或异常入口。 & \pass \\
\hline

42 & 用户端路由可访问。 & 访问 \url{/detect/paper} 等旧路径。 & 旧路径自动重定向至 \url{/upload}。 & \pass \\
\hline

43 & 管理端路由可访问。 & 访问 \url{/paper-logs} 等旧路径。 & 旧路径自动重定向至 \url{/logs}。 & \pass \\
\hline

\end{longtable}

\section{任务二（脚本测试）}

以下测试用例用于后端接口自动化测试。任务二共30项接口测试用例，覆盖认证与角色权限、文件上传与检测参数、人工审核流程约束、Review文本边界、论文与报告边界五类场景。接口基址为 \url{http://127.0.0.1:8000}；需认证接口使用Header：\texttt{Authorization: Bearer \textless token\textgreater}。根据最终测试结果，任务二全部接口测试用例均已通过。

\subsection{认证与角色权限测试}

\begin{longtable}{|P{0.8cm}|P{4.2cm}|P{4.0cm}|P{1.4cm}|P{4.3cm}|P{1.0cm}|}
\hline
\textbf{序号} & \textbf{测试用例说明} & \textbf{期望结果} & \textbf{实际状态码} & \textbf{Msg} & \textbf{结果} \\
\hline
\endfirsthead

\hline
\textbf{序号} & \textbf{测试用例说明} & \textbf{期望结果} & \textbf{实际状态码} & \textbf{Msg} & \textbf{结果} \\
\hline
\endhead

1 & 未携带Token访问 \url{/api/user/details/}。 & 接口拒绝访问，提示认证凭证缺失或无效。 & 401 & Authentication credentials were not provided & \pass \\
\hline

2 & 使用错误密码调用 \url{/api/login/} 登录普通用户。 & 登录失败，不返回access token。 & 401 & Invalid credentials & \pass \\
\hline

3 & 使用reviewer账号访问编辑侧人工审核申请创建接口 \url{/api/manual-review-requests/}。 & 接口拒绝非发布者发起人工审核申请。 & 403 & Only publishers can create... & \pass \\
\hline

4 & 使用publisher账号访问专家任务池 \url{/api/get_reviewer_tasks/}。 & 接口拒绝非专家查看专家任务池。 & 403 & Only reviewers can view tasks & \pass \\
\hline

5 & 使用普通用户账号调用 \url{/api/admin-login/}。 & 管理员登录失败，提示角色不匹配或无管理员权限。 & 400 & Invalid role. User is not an admin & \pass \\
\hline

6 & 用户A查询用户B的检测任务状态。 & 接口应拒绝越权访问或返回任务不存在。 & 404 & Detection task not found & \pass \\
\hline

\end{longtable}

\subsection{文件上传与检测参数校验测试}

\begin{longtable}{|P{0.8cm}|P{4.2cm}|P{4.0cm}|P{1.4cm}|P{4.3cm}|P{1.0cm}|}
\hline
\textbf{序号} & \textbf{测试用例说明} & \textbf{期望结果} & \textbf{实际状态码} & \textbf{Msg} & \textbf{结果} \\
\hline
\endfirsthead

\hline
\textbf{序号} & \textbf{测试用例说明} & \textbf{期望结果} & \textbf{实际状态码} & \textbf{Msg} & \textbf{结果} \\
\hline
\endhead

1 & 调用 \url{/api/upload/} 时不传 \texttt{file} 字段。 & 上传失败，提示缺少上传文件。 & 400 & No file provided & \pass \\
\hline

2 & 调用 \url{/api/upload/} 上传不支持的 \texttt{.exe} 文件。 & 上传失败，提示文件类型不合法。 & 400 & Unsupported upload type & \pass \\
\hline

3 & 调用 \url{/api/paper/upload/} 上传非PDF文件。 & 论文上传失败，提示论文检测仅支持PDF。 & 400 & 仅支持PDF文件 & \pass \\
\hline

4 & 调用 \url{/api/review/submit/} 时Review文本为空且未上传 \texttt{.txt} 文件。 & Review检测提交失败，提示输入内容不能为空。 & 400 & 请提供 text 或 txt 文件 & \pass \\
\hline

5 & 调用 \url{/api/detection/submit/} 时不传 \texttt{image\_ids} 或 \texttt{image\_ids} 为空数组。 & 检测任务创建失败，提示未提供图像ID。 & 400 & No image IDs provided & \pass \\
\hline

6 & 调用 \url{/api/detection/submit/} 提交不属于当前用户或当前组织的 \texttt{image\_id}。 & 检测任务创建失败，提示无权限提交或资源不存在。 & 404 & No valid images found & \pass \\
\hline

\end{longtable}

\subsection{人工审核流程约束测试}

\begin{longtable}{|P{0.8cm}|P{4.2cm}|P{4.0cm}|P{1.4cm}|P{4.3cm}|P{1.0cm}|}
\hline
\textbf{序号} & \textbf{测试用例说明} & \textbf{期望结果} & \textbf{实际状态码} & \textbf{Msg} & \textbf{结果} \\
\hline
\endfirsthead

\hline
\textbf{序号} & \textbf{测试用例说明} & \textbf{期望结果} & \textbf{实际状态码} & \textbf{Msg} & \textbf{结果} \\
\hline
\endhead

1 & 对不存在的 \texttt{detection\_task\_id} 发起人工审核申请。 & 申请创建失败，提示检测任务不存在。 & 404 & task not found & \pass \\
\hline

2 & 对无访问权限的检测任务发起人工审核申请。 & 申请创建失败，提示无权限访问该检测任务。 & 403 & 无权访问该检测任务 & \pass \\
\hline

3 & 管理员调用 \url{/api/handle_reviewRequest/<id>/} 时缺少 \texttt{choice} 参数。 & 审批失败，提示参数非法或记录不存在。 & 404 & ReviewRequest not found & \pass \\
\hline

4 & 管理员拒绝人工审核申请但未填写拒绝原因。 & 审批失败，提示拒绝时必须填写原因。 & 400 & 拒绝申请时必须填写原因 & \pass \\
\hline

5 & 专家调用 \url{/api/post_review/<manual_review_id>/} 提交不属于自己的人工审核任务。 & 提交失败，提示无权限操作该人工审核任务。 & 403 & 该用户没有审核的权限 & \pass \\
\hline

6 & 编辑在专家未完成前请求人工审核汇总或报告。 & 接口返回当前汇总状态或待完成提示，不出现服务端异常。 & 200 & 含 \texttt{admin\_state} 等字段 & \pass \\
\hline

\end{longtable}

\subsection{Review文本长度边界测试}

\begin{longtable}{|P{0.8cm}|P{4.2cm}|P{4.0cm}|P{1.4cm}|P{4.3cm}|P{1.0cm}|}
\hline
\textbf{序号} & \textbf{测试用例} & \textbf{期望结果} & \textbf{实际状态码} & \textbf{Msg} & \textbf{结果} \\
\hline
\endfirsthead

\hline
\textbf{序号} & \textbf{测试用例} & \textbf{期望结果} & \textbf{实际状态码} & \textbf{Msg} & \textbf{结果} \\
\hline
\endhead

1 & 空字符串。 & 拒绝，提示Review内容不能为空。 & 400 & 请提供 text 或 txt 文件 & \pass \\
\hline

2 & 仅包含空格、换行或制表符。 & 拒绝，清洗后为空文本。 & 400 & 请提供 text 或 txt 文件 & \pass \\
\hline

3 & 极短文本，例如“Hi”。 & 接受或给出最小长度提示，不应出现服务端异常。 & 200 & \texttt{task\_id} 返回 & \pass \\
\hline

4 & 正常长度Review文本，约数百字。 & 接受，返回任务编号、状态和清洗后文本长度。 & 200 & 含 \texttt{cleaned\_text\_length} & \pass \\
\hline

5 & 超长Review文本，约20万字。 & 拒绝或截断处理并给出明确提示，不应导致超时或500错误。 & 400 & 超过最大长度100000字符 & \pass \\
\hline

6 & 包含BOM、控制字符、连续空白和混合编码字符。 & 接受并完成文本清洗，长度与清洗后文本一致。 & 200 & \texttt{cleaned\_text\_length=44} & \pass \\
\hline

\end{longtable}

\subsection{论文与报告边界测试}

\begin{longtable}{|P{0.8cm}|P{4.2cm}|P{4.0cm}|P{1.4cm}|P{4.3cm}|P{1.0cm}|}
\hline
\textbf{序号} & \textbf{测试用例} & \textbf{期望结果} & \textbf{实际状态码} & \textbf{Msg} & \textbf{结果} \\
\hline
\endfirsthead

\hline
\textbf{序号} & \textbf{测试用例} & \textbf{期望结果} & \textbf{实际状态码} & \textbf{Msg} & \textbf{结果} \\
\hline
\endhead

1 & 上传0KB空PDF。 & 拒绝，提示文件为空或解析失败。 & 400 & failed to open PDF & \pass \\
\hline

2 & 上传正常可解析PDF。 & 接受，返回论文文件编号、文件名、上传时间、段落数量等信息。 & 200 & \texttt{paper\_file\_id}、\texttt{paragraph\_count} 等 & \pass \\
\hline

3 & 上传损坏或加密PDF。 & 拒绝，提示PDF解析失败或文件不合法。 & 400 & 解析失败或 encrypted & \pass \\
\hline

4 & 查询不存在的 \texttt{task\_id} 状态。 & 拒绝或返回任务不存在提示。 & 404 & 任务不存在 & \pass \\
\hline

5 & 下载未完成检测任务的报告。 & 拒绝，提示任务尚未完成或报告未生成。 & 400 & Task not completed yet & \pass \\
\hline

6 & 下载已完成检测任务的报告。 & 接受，返回PDF文件流。 & 200 & PDF文件流 & \pass \\
\hline

\end{longtable}

\section*{结论}

本轮测试覆盖用户端、管理端与后端REST API的核心业务路径。任务一共43项功能测试用例，任务二共30项接口测试用例，合计73项测试用例均已执行并通过。系统能够支持编辑发起检测、查看历史与报告、申请人工审核，支持组织管理员审批人工审核申请，支持专家提交人工审核结论，并支持管理端查看检测任务、资源、用户组织与系统日志。测试过程中发现的问题均已完成修复并复测通过。

\section*{遗留说明}

本轮测试未纳入性能压测、WebSocket实时推送定量验收、深度跨模态融合算法效果验证以及精准模式LLM效果差异评估。当前论文检测仅支持PDF格式，DOCX文件在批量入口提示转换为PDF后再提交。

\end{document}